% page 95 du fac-simile (image p-094 du scan)
% bloc 170.2 x 229.0 mm ; marge gauche 31.8 mm
\pgc{10.160mm}{0.000mm}{
\l{\cel{56}87}
\l{}
\l{}
\l{\sou{chapleto.}\cel{2}Asemblajo de kin deki de grani filizita, separita per}
\l{\cel{3}grano plu grosa. - EFL.}
\l{}
\l{\sou{charo}. I. Veturo lejera di la antiqui, kun du roti, aperturoza}
\l{\cel{3}dope, tirata da du o plura kavali jungita tale ke li avancas}
\l{\cel{3}fronte, e sur la parto avana di qua la duktanto stacas. -}
\l{\cel{3}II. Veturo dekorita, en la festi, procesioni, masko-festi,}
\l{\cel{3}e qua portas personegi simbola, historiala, maskiti, e c.-DEFIS.}
\l{}
\l{\sou{charioto}. Kargo-veturo kun quar roti, rideloza ed un-timona,}
\l{\cel{3}uzata en la ruro por transportar la rekoltaji, la vari vend-}
\l{\cel{3}enda, e c. - en la armei, por transportar la bagaji, e c. - F.}
\l{}
\l{\sou{charjar}. (\sou{trans}.) I. Lokizar (veturo, e c.) sub la pezo di la}
\l{\cel{3}objekti (kargajo) portenda ulube, transportenda. - II. Garnisar}
\l{\cel{3}(ulo) ye pezajo, ye ula quanto, destinita a ta o ca uzo (fusilo,}
\l{\cel{3}pipo, bobino, e c.) - EFIS.}
\l{}
\l{\sou{charmar}.\cel{2}(\sou{trans}.)Ravisar, seduktar per atraktivo forta. - EF.}
\l{}
\l{\sou{charniro}. Ligilo ek du fero-plaketi, artikoza, de qui adminime}
\l{\cel{3}un povas movar cirkum axo. - EFIES.}
\l{}
\l{\sou{charto}. En la mezepoko, omna akto en qua enrejistresis la akti}
\l{\cel{3}di proprieto, vendo, privileji grantita, e c. - DEFIRS.}
\l{}
\l{\sou{chasar}. (\sou{trans.}) Proxime-sequar la animali por kaptar o destruktar}
\l{\cel{3}li. - EFIS.}
\l{}
\l{\sou{chasta}. Qua konservas sexuo-pura sua korpo e sua psiko. - EFIS.}
\l{}
\l{\sou{chazublo}. Vesto sacerdotala per qua kovresas la sacerdoto dum ke}
\l{\cel{3}lu celebras la santa sakrifiko, qua pendas sur la albo, ed havas}
\l{\cel{3}kom ornivo precipua la figuro di la Kruco (kristana). - EF.}
\l{}
\l{\sou{che}.\cel{2}En la domo, habiteyo, lando, domeno (materiala o spiritala}
\l{\cel{3}- inter ta o ca sorto de grupani. - F.}
\l{}
\l{\sou{chefo}. La persono qua esas ye la kapo di grupo quan lu duktas ed}
\l{\cel{3}a qua lu imperas. - DEFRS.}
\l{}
\l{\sou{chelato}. Kasko di qua la kaloto esas mi-sferatra, la viziero kurta}
\l{\cel{3}e fixa, la nukumo granda, quan (en la mezepoko ed en la yarcento}
\l{\cel{3}XVI-esma) +weris la militisti kavaliera. - EFIS.}
\l{}
\l{\sou{cheko}. (\sou{Banko}).Folieto separita de kayero talonoza, kompozita ek}
\l{\cel{3}bilieti pagebla ye vido (da banko) e qua, pro pagendeso a la}
\l{\cel{3}portanto, esas indosebla. - DEFRS.}
\l{}
\l{\sou{chera}. Qua kustas plu multe kam lo analoga. - FIS.}
}
